\section{Related Work}
\label{sec:related}

\subsection{Multi-Task Model Merging}
Model merging has emerged as a crucial approach to multi-task learning without retraining. Early methods explored simple parameter averaging across model checkpoints fine-tuned on the same data, known as Model Soups \cite{wortsman2022model}. For diverse downstream tasks, Task Arithmetic \cite{ilharco2023editing} demonstrated that subtracting pre-trained weights from fine-tuned expert weights yields ``task vectors'' that can be linearly added to the base weights to transfer multiple capabilities. 

However, linear combination suffers from parameter interference and task dominance. To address this, RegMean \cite{jin2023regmean} utilizes linear regression over layer activations to compute closed-form merging matrices, but requires computing activation covariance matrices on calibration data. Fisher Weighted Averaging \cite{matena2021fisher, donyehiya2023fisher} uses diagonal Fisher information as parameter importance weights, which is computationally expensive and data-dependent. TIES-Merging \cite{yadav2023resolving} resolves sign agreements and prunes task vectors by a set threshold, followed by scaling. While effective, TIES introduces several complex heuristic steps and multiple sensitive hyperparameters (e.g., pruning ratios, scaling factors). In contrast, NETA resolves representation dominance analytically in a single step with zero additional hyperparameters, preserving all task updates without heuristic pruning.

\subsection{Test-Time Weight Adaptation}
To dynamically scale the contributions of different tasks, test-time adaptation methods have been proposed. AdaMerging \cite{yang2024adamerging} optimizes layer-wise and task-wise coefficients on a small, unlabeled calibration set using gradient-based minimization of joint prediction entropy. SyMerge \cite{yu2023symerge} integrates symmetry-aware parameter correspondence into weight merging to further refine performance. Sharpness-Aware Isotropic Merging (SAIM) \cite{jordon2024saim} attempts to locate flatter minima in the weight space through global perturbations during merging. 

While test-time adaptation often improves metrics, it exhibits severe drawbacks: it is computationally intensive, requires optimization loops, and introduces hyperparameters (learning rate, epochs, and schedules). More critically, as we demonstrate in this work, unconstrained optimization of multi-task weight coefficients via joint entropy minimization is highly susceptible to transductive overfitting and task bias (the Overfitting-Optimizer Paradox). NETA completely avoids these limitations by using a training-free, parameter-free closed-form transformation that operates zero-shot directly on the expert weights.

\subsection{Loss Landscapes and Mode Connectivity}
Linear model merging relies on the unique geometry of deep network loss landscapes. Modern networks initialized from a shared pre-trained checkpoint and independently fine-tuned have been shown to reside in a shared, low-loss basin \cite{neyshabur2020what, wortsman2022model}. Under this shared basin condition, the fine-tuned expert models are connected by low-loss paths \cite{garipov2018mode, draxler2018essentially}, making direct linear weight interpolation viable without devastating representation collapse. Git Re-Basin \cite{ainsworth2023git} and ZipIt! \cite{stoica2023zipit} expand on this by matching permutation symmetries between models. NETA builds on the shared-basin assumption of Task Arithmetic, showing that within this low-loss manifold, we can perform isotropic balancing of task update scales directly in the parameter space to achieve optimal multi-task representation alignment.
